% titel: Blatt 6, A2: Lewis-Strukturformeln zeichnen
% skript: Mortimer Kap. 8 (chemie: Atombindung, Lewis-Formeln)

Eine Lewis-Struktur zeigt alle Valenzelektronen als Bindungspaare (Striche) und
freie Elektronenpaare (Punkte). Man zeichnet sie nach einem festen Schema.

\begin{werkzeugkasten}[Lewis-Formel nach Schema]
\begin{enumerate}
  \item Valenzelektronen aller Atome zählen; bei Ionen Ladung mitrechnen
        ($+$ aufläden: Elektron entfernen, $-$: Elektron addieren).
  \item Das am wenigsten elektronegative Atom zum Zentrum machen (ausser H).
  \item Atome über Einfachbindungen verbinden, dann verbleibende Elektronen als
        freie Paare auf die Peripherie verteilen --- oft reicht das.
  \item        Reicht das Oktett nicht, Doppel-/Dreifachbindungen einbauen (zentrale
        Elemente der 2.\ Periode wie $\mathrm{C,N,O}$ dürfen nur maximal ein Oktett).
  \item Für zentrale Elemente ab Periode 3 ist Oktetterweiterung $>8$ erlaubt.
\end{enumerate}
\textit{Skript:} Mortimer Kap. 8 (Lewis-Formeln, Oktettregel).
\end{werkzeugkasten}

\begin{achtung}[Darauf achten]
\begin{itemize}
  \item Formalladungen: prüfen, dass die Summe der Formalladungen der
        Ionenladung entspricht (z.\ B. $\mathrm{-}$ bei $\mathrm{PCl_6^-}$).
  \item H und He kommen nie ins Zentrum und bilden nur Einfachbindungen.
  \item Bei Molekülen mit $\mathrm{SO}$- oder $\mathrm{NO}$-Teil struktur die
        Mehrfachbindung so wählen, dass möglichst alle Oktetts erfüllt sind
        (formal stabilste Struktur).
  \item Bei Periode-2-Zentren gibt es keine Oktetterweiterung --- das zwingt z.B.
        in $\mathrm{CO}$ zu einer Dreifachbindung.
\end{itemize}
\end{achtung}
